% Provenance (never printed): see sources/outline_sources.md (Year 2: radiative forcing
% quantified; feedbacks and climate sensitivity; detection and attribution; scenarios; the
% remaining carbon budget; components of sea-level rise; physical science only). Every statement
% from IPCC AR6 WG1 (SPM; chapters 3, 5, 7 and 9; Annex III data), ledger rows in
% sources/ledger/book2.md (ipcc:*) and sources/ledger/book4.md (ar6*:*). Curves and numbers from
% figdata/bachelor-2/anthropogenic.py and the Book 4 second-writer scratch calculator calc/ch27.py.
\chapter{El cambio climático antropógeno}\label{ch:b2:anthropogenic-climate-change}

Desde 1970, la superficie del globo se ha calentado más deprisa que en cualquier otro periodo de cincuenta años de, al menos,
los últimos dos mil años (\emph{confianza alta}). Los volúmenes escolares contaron qué ha cambiado y
por qué. Este capítulo convierte ese relato en un balance: cuántos vatios por metro cuadrado añaden los
factores humanos, cuántos devuelve al espacio el planeta más caliente, cuánto calentamiento produce cada tonelada de dióxido
de carbono y cómo se reparte la subida del mar entre el agua que se calienta y el hielo que se funde. Cada cifra
viene con su incertidumbre, en el lenguaje calibrado de las evaluaciones.
% ledger: ipcc:fastest-2000yr, ipcc:likelihood

\begin{recall}
\textbf{De los volúmenes escolares}: el forzamiento radiativo, $\qty{3.93}{W.m^{-2}}$ por una duplicación del
dióxido de carbono, y su ley logarítmica; la retroalimentación climática; la \omterm{def:b2:anthropogenic-climate-change:ecs}{sensibilidad climática de equilibrio}, con una mejor
estimación de \qty{3}{\celsius}, \emph{probable} entre 2.5 y \qty{4}{\celsius}; la \omterm{def:b2:anthropogenic-climate-change:da}{atribución}, el escenario de emisiones
y la proyección; el lenguaje calibrado: \emph{probable}, de 66 a \qty{100}{\percent}; \emph{muy
probable}, de 90 a \qty{100}{\percent}; un intervalo $[x \text{ a } y]$ es un intervalo muy probable. La influencia
humana ha calentado el clima en \qty{1.07}{\celsius} entre 1850--1900 y 2010--2019.
\textbf{De los libros de física}: el cuerpo negro y la ley de Stefan--Boltzmann; la dilatación térmica.
\textbf{De los libros de biología}: el ciclo del carbono, la fracción que permanece en el aire, los sumideros de carbono y el
balance de carbono de la biosfera. \textbf{Del \cref{ch:b2:deep-time-climate}}: el ciclo lento del carbono,
que solo retira un exceso de dióxido de carbono a lo largo de cientos de miles de años.
\end{recall}
% ledger: ipcc:erf-2xco2, ipcc:log-co2, ipcc:ecs, ipcc:likelihood, ipcc:attribution

\section{El forzamiento cuantificado}\label{sec:b2:anthropogenic-climate-change:forcing}

\begin{definition}[Forzamiento radiativo efectivo]\label{def:b2:anthropogenic-climate-change:erf}
El \emph{forzamiento radiativo efectivo}\index{forzamiento radiativo efectivo} (ERF, por sus siglas en inglés) de una perturbación
---un cambio en un gas de efecto invernadero, en los aerosoles, en el Sol--- es el cambio que provoca en el flujo neto
descendente de radiación en el techo de la atmósfera, en \unit{W.m^{-2}}, una vez que la atmósfera,
las nubes y la superficie terrestre se han ajustado a ella, pero antes de que la temperatura global de la superficie
haya respondido. Los ajustes que se cuentan son los causados por el propio agente, con independencia de la
temperatura de la superficie.
\end{definition}
% ledger: ar6ch7:erf-def

La definición traza la línea entre forzamiento y respuesta según lo que depende de la temperatura de la
superficie, no según la rapidez. El forzamiento instantáneo del dióxido de carbono es el cambio de flujo justo
después de añadir el gas; la estratosfera, que se enfría en unos meses, lo modifica una primera vez;
la troposfera, su vapor de agua y sus nubes lo modifican de nuevo. Para una duplicación del dióxido de
carbono, la evaluación halla \qty{3.75}{W.m^{-2}} una vez ajustada la estratosfera, y unos
ajustes en la troposfera que se compensan en parte: \qty{-0.60}{W.m^{-2}} por su calentamiento, +0.22
por el vapor de agua, +0.45 por las nubes y +0.11 por la superficie terrestre, \qty{+0.18}{W.m^{-2}} en total.
El ERF es $3.75 + 0.18 = \qty{3.93}{W.m^{-2}}$, un \qty{6}{\percent} más que los
\qty{3.71}{W.m^{-2}} de la evaluación anterior (\cref{exo:b2:anthropogenic-climate-change:1}).
% ledger: ar6ch7:2xco2, ar6ch7:erf-def

El forzamiento de cada agente desde 1750 se ha reconstruido año a año. La figura siguiente lo
muestra: el dióxido de carbono, con \qty{2.16}{W.m^{-2}} en 2019; los demás gases de efecto invernadero ---metano,
óxido nitroso, gases halogenados y ozono---, con \qty{1.69}{W.m^{-2}} en conjunto; y los
aerosoles, con \qty{-1.06}{W.m^{-2}}. El forzamiento humano total es de \qty{2.72}{W.m^{-2}}, y su intervalo muy
probable, de 1.96 a \qty{3.48}{W.m^{-2}}, es amplio sobre todo por los aerosoles. Su ERF en
1750--2014 se evalúa en $-1.3$ [$-2.0$ a $-0.6$] \unit{W.m^{-2}} (\emph{confianza media}),
la mayor parte a través de su efecto sobre las nubes, que se vuelven más brillantes cuando los aerosoles aportan más gotitas.
Las erupciones volcánicas añaden picos negativos breves (\cref{ch:b2:eruptive-styles}): el más profundo del
registro, \qty{-4.8}{W.m^{-2}} en 1816, sigue a la erupción del Tambora de 1815.
% ledger: data:b4-ar6-erf, ipcc:erf-total, ipcc:erf-agents, ar6ch7:aerosol, usgs:msh-comparisons

\begin{omfigure}
\centering
\begin{tikzpicture}
\begin{axis}[omrecord, width=0.92\linewidth, height=6.2cm, xmin=1750, xmax=2020, ymin=-5.2,
  ymax=3.8, xlabel={año}, ylabel={ERF desde 1750 (\unit{W.m^{-2}})},
  xtick={1750,1800,...,2000}, /pgf/number format/1000 sep={}, xtick pos=bottom, ytick pos=left,
  tick label style={font=\scriptsize}, label style={font=\small},
  legend style={at={(0.5,-0.24)}, anchor=north, font=\scriptsize, draw=black!30,
  legend columns=3, /tikz/every even column/.append style={column sep=6pt}},
  legend cell align=left]
% computed in figdata/bachelor-2/anthropogenic.py part erf; ledger: data:b4-ar6-erf
\addplot[name path=hi, draw=none, forget plot] table[x=year, y=hi] {figdata/out/bachelor-2-anthropogenic-erf.dat};
\addplot[name path=lo, draw=none, forget plot] table[x=year, y=lo] {figdata/out/bachelor-2-anthropogenic-erf.dat};
\addplot[black!15, forget plot] fill between[of=hi and lo];
\addplot[black!45, line width=0.4pt] table[x=year, y=natural] {figdata/out/bachelor-2-anthropogenic-erf.dat};
\addlegendentry{natural (volcanes, Sol)}
\addplot[omThm, line width=1pt] table[x=year, y=co2] {figdata/out/bachelor-2-anthropogenic-erf.dat};
\addlegendentry{dióxido de carbono}
\addplot[omThm!50!orange, line width=1pt, densely dashed] table[x=year, y=other_ghg] {figdata/out/bachelor-2-anthropogenic-erf.dat};
\addlegendentry{otros gases de efecto invernadero y ozono}
\addplot[omDef, line width=1pt] table[x=year, y=aerosol] {figdata/out/bachelor-2-anthropogenic-erf.dat};
\addlegendentry{aerosoles}
\addplot[black, line width=1.4pt] table[x=year, y=anthro] {figdata/out/bachelor-2-anthropogenic-erf.dat};
\addlegendentry{total humano (gris: 5--95\,\%)}
\draw[black!40] (axis cs:1750,0) -- (axis cs:2020,0);
\node[anchor=west, font=\scriptsize, text=black!60, inner sep=1pt] at (axis cs:1817.5,-4.75) {$\leftarrow$ Tambora};
\end{axis}
\end{tikzpicture}
\par\smallskip{\small \omterm{def:b2:anthropogenic-climate-change:erf}{Forzamiento radiativo efectivo} de 1750 a 2019, respecto de 1750, por factor.
El forzamiento humano total alcanza \qty{2.72}{W.m^{-2}} en 2019, con el intervalo muy probable sombreado.
Datos: IPCC, Sexto Informe de Evaluación, Grupo de Trabajo~I, datos ampliados del Anexo~III (CC BY 4.0).}
\end{omfigure}
% values: anthro 2.7204 (1.9569-3.4819), co2 2.1563, other_ghg 1.6850, aerosol -1.0583, volcanic -4.8162 in 1816 (anthropogenic.py numbers)

\section{El balance de energía}\label{sec:b2:anthropogenic-climate-change:energy}

Un forzamiento positivo hace que el planeta gane energía. A medida que se calienta, radia más hacia el espacio, y la
ganancia disminuye; el clima solo se estabiliza cuando ambos se equilibran.

\begin{definition}[Desequilibrio energético de la Tierra; parámetro de retroalimentación climática]\label{def:b2:anthropogenic-climate-change:eei}
El \emph{desequilibrio energético de la Tierra}\index{desequilibrio energetico de la Tierra@desequilibrio energético de la Tierra} $N$ es el flujo neto de energía que entra en
el sistema climático por el techo de la atmósfera, promediado sobre el globo y sobre un año o más,
en \unit{W.m^{-2}}: la luz solar absorbida menos el infrarrojo saliente. Es el ritmo al que el océano,
la tierra, el hielo y el aire ganan calor, por metro cuadrado de la superficie terrestre.
\end{definition}

\begin{definition}[Parámetro de retroalimentación climática]\label{def:b2:anthropogenic-climate-change:lambda}
El \emph{parámetro de retroalimentación climática}\index{parametro de retroalimentacion climatica@parámetro de retroalimentación climática} $\lambda$ es el cambio del
flujo neto descendente en el techo de la atmósfera por grado de calentamiento global de la superficie, en
\unit{W.m^{-2}.\celsius^{-1}}, causado por la respuesta del sistema climático a ese calentamiento. Un
clima estable tiene $\lambda < 0$: el calentamiento aumenta la pérdida de energía hacia el espacio.
\end{definition}

El volumen escolar escribía el balance con el módulo de $\lambda$, un número positivo; a partir de aquí se
conserva el signo, como en las evaluaciones, que llaman $\alpha$ a este parámetro.
% ledger: ar6ch7:energy-budget

\begin{proposition}[Balance global de energía]\label{prop:b2:anthropogenic-climate-change:balance}
Sean $\Delta F$ el \omterm{def:b2:anthropogenic-climate-change:erf}{forzamiento radiativo efectivo} y $\Delta T$ el calentamiento global de la superficie
respecto de un estado de equilibrio. Si la respuesta del flujo saliente es lineal en $\Delta T$,
\[
  \Delta N = \Delta F + \lambda\,\Delta T .
\]
En un nuevo equilibrio $\Delta N = 0$, así que $\Delta T_{\text{eq}} = -\Delta F/\lambda$.
\end{proposition}
\begin{proof}
Por definición, el forzamiento cambia el flujo neto descendente en $\Delta F$ a temperatura de la superficie
fija, y la respuesta al calentamiento lo cambia en $\lambda\,\Delta T$ a primer orden;
el desequilibrio es su suma. Igualarlo a cero da el equilibrio.
\end{proof}

Medir $N$ significa medir el calor que gana el sistema, y el océano almacena casi todo ese calor.
De 1971 a 2018, el sistema climático ganó \qty{435}{ZJ} (zettajulios, \qty{e21}{J}): el océano, un
\qty{91}{\percent}, \qty{242}{ZJ} en sus \qty{700}{m} superiores y \qty{154}{ZJ} por debajo; la tierra, un
\qty{5}{\percent}; el hielo que se funde, un \qty{3}{\percent}; la atmósfera, un \qty{1}{\percent}. El
desequilibrio medio fue de \qty{0.57}{W.m^{-2}} en 1971--2018, y ha crecido: \qty{0.79}{W.m^{-2}}
[0.52 a 1.06] en 2006--2018 (\cref{exo:b2:anthropogenic-climate-change:2}).
% ledger: ar6ch7:inventory, ipcc:heating

\begin{omfigure}
\centering
\begin{tikzpicture}[font=\small, >=Stealth]
\draw[fill=omDef!8, draw=omDef!60] (-3.2,0) rectangle (3.2,1.6);
\node[align=center] at (0,1.05) {sistema climático};
\node[align=center, font=\scriptsize] at (0,0.5) {océano \qty{91}{\percent}, tierra \qty{5}{\percent}, hielo \qty{3}{\percent}, aire \qty{1}{\percent}};
\draw[->, line width=2.2pt, omThm] (-2.0,3.3) -- node[pos=0.12, left=3pt, align=right, font=\scriptsize] {forzamiento $\Delta F$} (-2.0,1.7);
\draw[->, line width=2.2pt, omDef] (2.0,1.7) -- node[pos=0.88, right=3pt, align=left, font=\scriptsize] {respuesta $\lambda\,\Delta T$} (2.0,3.3);
\draw[->, line width=1.4pt, black!70] (0,-0.1) -- node[right=3pt, align=left, font=\scriptsize] {almacenado: $\Delta N = \Delta F + \lambda\,\Delta T$} (0,-1.0);
\draw[black!40, dashed] (-3.4,2.4) -- (3.4,2.4) node[right, font=\scriptsize, text=black!60] {techo de la atmósfera};
\end{tikzpicture}
\par\smallskip{\small El balance de energía de un clima forzado (esquema). El forzamiento añade energía en
el techo de la atmósfera, la respuesta del planeta más caliente retira una parte ($\lambda\,\Delta T$ es
negativo) y la diferencia se almacena, sobre todo en el océano. Reparto del calor ganado de 1971
a 2018.}
\end{omfigure}
% ledger: ar6ch7:inventory

\begin{method}[Una sensibilidad climática a partir del balance de energía histórico]\label{met:b2:anthropogenic-climate-change:energy-budget}
\begin{enumerate}
  \item Elíjanse dos periodos, uno de referencia cercano al equilibrio y otro reciente; tómense el cambio del
        forzamiento $\Delta F$, del desequilibrio $\Delta N$ y de la temperatura global $\Delta T$.
  \item Despéjese del balance el parámetro de retroalimentación: $\lambda = (\Delta N - \Delta F)/\Delta T$.
  \item La sensibilidad a una duplicación es entonces $-F_{2\times}/\lambda$; el calentamiento en el momento del
        forzamiento, escalado a $F_{2\times}$, es $F_{2\times}\Delta T/\Delta F$.
\end{enumerate}
\end{method}

De 1850--1900 a 2006--2019, la evaluación toma $\Delta N = \qty{0.59}{W.m^{-2}}$, $\Delta T =
\qty{1.03}{\celsius}$ y $\Delta F = \qty{2.20}{W.m^{-2}}$. Entonces $\lambda = (0.59 - 2.20)/1.03 =
\qty{-1.56}{W.m^{-2}.\celsius^{-1}}$, y la sensibilidad es $3.93/1.56 = \qty{2.5}{\celsius}$,
el valor que halla la evaluación, con un intervalo muy probable de 1.6 a \qty{4.8}{\celsius}. Pero se trata de
una sensibilidad \emph{efectiva}, para las retroalimentaciones del calentamiento ocurrido hasta ahora. El calentamiento ha sido
más intenso donde las retroalimentaciones son más negativas, en el Pacífico tropical occidental, y más débil en el
Pacífico oriental y el océano Austral; a medida que esa distribución evolucione, se espera que el parámetro de retroalimentación
se vuelva menos negativo (\emph{confianza alta}), de modo que la sensibilidad a largo plazo es mayor que la
estimación histórica (\cref{exo:b2:anthropogenic-climate-change:7}).
% ledger: ar6ch7:energy-budget, ar6ch7:pattern, ipcc:erf-2xco2
% values: lambda_energy -1.563, ecs_energy 2.514 (anthropogenic.py numbers)

\section{Retroalimentaciones y sensibilidad}\label{sec:b2:anthropogenic-climate-change:feedbacks}

\begin{definition}[Respuesta de Planck; retroalimentaciones del vapor de agua, del gradiente vertical y de las nubes]\label{def:b2:anthropogenic-climate-change:feedbacks}
El parámetro de retroalimentación se descompone en las contribuciones de las magnitudes que cambian con la
temperatura de la superficie. La \emph{respuesta de Planck}\index{respuesta de Planck} es el infrarrojo adicional
emitido hacia el espacio cuando la superficie y toda la troposfera se calientan en la misma cantidad. La
\emph{retroalimentación del vapor de agua}\index{retroalimentacion del vapor de agua@retroalimentación del vapor de agua} es el cambio de flujo debido al aumento
del vapor de agua atmosférico con la temperatura. La \emph{retroalimentación del gradiente vertical de temperatura}
\index{retroalimentacion del gradiente vertical de temperatura@retroalimentación del gradiente vertical de temperatura} es el cambio de flujo debido a un calentamiento que no es uniforme con la altura. La \emph{retroalimentación de las
nubes}\index{retroalimentacion de las nubes@retroalimentación de las nubes} es el cambio de flujo debido al cambio de la cantidad, la altura y el
brillo de las nubes.
\end{definition}

\begin{proposition}[La respuesta de Planck]\label{prop:b2:anthropogenic-climate-change:planck}
Si la Tierra emite hacia el espacio una fracción fija $1 - g$ del infrarrojo $\sigma T^4$ de su superficie a la
temperatura $T$, su \omterm{def:b2:anthropogenic-climate-change:feedbacks}{respuesta de Planck} es
\[
  \lambda_{\text P} = -(1 - g)\,4\sigma T^3 .
\]
Con $g = 0.4$ y $T = \qty{288}{K}$, $\lambda_{\text P} = \qty{-3.25}{W.m^{-2}.K^{-1}}$.
\end{proposition}
\begin{proof}
El flujo saliente es $(1 - g)\sigma T^4$; su derivada respecto de $T$ es $(1 - g)\,4\sigma
T^3$, y una pérdida adicional es una contribución negativa al flujo neto descendente. Numéricamente, $0.6
\times 4 \times \num{5.670e-8} \times 288^3 = 3.25$.
\end{proof}
% ledger: ar6ch7:planck, const:sigma, ipcc:budget-surface
% values: planck -3.2509 (anthropogenic.py numbers)

La fracción $g$ es el efecto invernadero normalizado: la atmósfera retiene \qty{159}{W.m^{-2}} de
los \qty{398}{W.m^{-2}} que emite la superficie, $g \approx 0.4$. La evaluación, con códigos de radiación
detallados, halla $-3.22$ [$-3.4$ a $-3.0$] (\emph{confianza alta}). El vapor de agua sigue
la ley de Clausius--Clapeyron a escala global, y más vapor absorbe más infrarrojo: la
\omterm{def:b2:anthropogenic-climate-change:feedbacks}{retroalimentación del vapor de agua} es fuertemente positiva, \qty{1.85}{W.m^{-2}.\celsius^{-1}} según los satélites.
En los trópicos, la troposfera superior se calienta más que la superficie y radia entonces más: la
retroalimentación del gradiente vertical es negativa allí, positiva cerca de los polos y de unos $-0.50$ a escala global. Las dos
se evalúan juntas, \num{1.30} [1.1 a 1.5], porque lo que un modelo pone en una tiende a
quitarlo de la otra. La nieve y el hielo que se funden oscurecen la superficie, +0.35; las nubes añaden \num{0.42}
[$-0.10$ a $0.94$], la mayor incertidumbre. Una retroalimentación neta negativa de las nubes es \emph{muy improbable}.
% ledger: ar6ch7:planck, ar6ch7:wv-lr, ar6ch7:table710, ipcc:cloud-negative, ipcc:budget-surface

\begin{omfigure}
\centering
\begin{tikzpicture}
\begin{axis}[xbar, width=0.78\linewidth, height=5.4cm, xmin=-3.7, xmax=1.8, ymin=0.4, ymax=6.6,
  ytick={1,...,6}, yticklabels={neta, biogeofísicas y químicas, nubes, albedo de la superficie,
  vapor de agua + gradiente, Planck}, y tick label style={font=\scriptsize},
  tick label style={font=\scriptsize}, xlabel={parámetro de retroalimentación (\unit{W.m^{-2}.\celsius^{-1}})},
  label style={font=\small}, bar width=7pt, axis x line=bottom, axis y line=left,
  axis line style={-}, xmajorgrids, grid style={gray!22, line width=0.3pt}, extra x ticks={0},
  extra x tick style={grid style={black!60, line width=0.5pt}}]
% ledger: ar6ch7:table710 (central values and very likely ranges)
\addplot[fill=omDef!55, draw=omDef!80!black, error bars/.cd, x dir=both, x explicit,
  error bar style={black!70}] coordinates {
  (-3.22,6) += (0.22,0) -= (0.18,0)
  (-0.01,2) += (0.26,0) -= (0.26,0)
  (-1.16,1) += (0.65,0) -= (0.65,0)};
\addplot[fill=omThm!55, draw=omThm!80!black, error bars/.cd, x dir=both, x explicit,
  error bar style={black!70}] coordinates {
  (1.30,5) += (0.20,0) -= (0.20,0)
  (0.35,4) += (0.25,0) -= (0.25,0)
  (0.42,3) += (0.52,0) -= (0.52,0)};
\end{axis}
\end{tikzpicture}
\par\smallskip{\small Las retroalimentaciones climáticas evaluadas por el IPCC (Sexto Informe de Evaluación, Grupo de
Trabajo~I, tabla 7.10), valores centrales e intervalos muy probables. Las retroalimentaciones positivas amplifican un calentamiento
(rojo), las negativas lo amortiguan (azul); la neta es la suma.}
\end{omfigure}

\begin{definition}[Sensibilidad climática de equilibrio]\label{def:b2:anthropogenic-climate-change:ecs}
La \emph{sensibilidad climática de equilibrio}\index{sensibilidad climatica de equilibrio@sensibilidad climática de equilibrio} (ECS, por sus siglas en inglés) es el
calentamiento global de la superficie, en el nuevo equilibrio, causado por una duplicación del dióxido de carbono del
aire desde su nivel preindustrial, con los casquetes de hielo mantenidos fijos.
\end{definition}

\begin{proposition}[Sensibilidad y retroalimentación]\label{prop:b2:anthropogenic-climate-change:ecs}
Si el parámetro de retroalimentación es constante, $\text{ECS} = -F_{2\times}/\lambda$. Con $F_{2\times} =
\qty{3.93}{W.m^{-2}}$ y el $\lambda = \qty{-1.16}{W.m^{-2}.\celsius^{-1}}$ neto de las retroalimentaciones
anteriores, $\text{ECS} = \qty{3.4}{\celsius}$; la \omterm{def:b2:anthropogenic-climate-change:feedbacks}{respuesta de Planck} sola daría $3.93/3.22 =
\qty{1.2}{\celsius}$.
\end{proposition}
\begin{proof}
Aplíquese la \cref{prop:b2:anthropogenic-climate-change:balance} en el equilibrio con $\Delta F =
F_{2\times}$.
\end{proof}
% ledger: ar6ch7:table710, ipcc:erf-2xco2
% values: ecs_process 3.388 (anthropogenic.py numbers)

Las retroalimentaciones son una línea de evidencia; el calentamiento desde 1850, el balance de energía y los climas
del pasado son otras. Combinadas, dan la ECS evaluada de \qty{3}{\celsius}, \emph{probable} entre
2.5 y \qty{4}{\celsius} (\emph{confianza alta}), \emph{muy probable} entre 2 y \qty{5}{\celsius}.
Como la ECS varía como $1/\lambda$, una incertidumbre simétrica sobre $\lambda$ da un intervalo
de ECS estirado hacia los valores altos (\cref{exo:b2:anthropogenic-climate-change:12}).
% ledger: ipcc:ecs

\begin{definition}[Respuesta climática transitoria]\label{def:b2:anthropogenic-climate-change:tcr}
La \emph{respuesta climática transitoria}\index{respuesta climatica transitoria@respuesta climática transitoria} (TCR, por sus siglas en inglés) es el calentamiento global de la superficie
en el momento en que el dióxido de carbono del aire se ha duplicado, cuando aumenta un \qty{1}{\percent}
al año desde su nivel preindustrial, unos setenta años después del comienzo.
\end{definition}

La TCR es menor que la ECS porque el océano sigue absorbiendo calor cuando se alcanza la duplicación.
Su mejor estimación es \qty{1.8}{\celsius}, \emph{probable} entre 1.4 y \qty{2.2}{\celsius}
(\emph{confianza alta}); el balance de energía histórico por sí solo da $3.93 \times 1.03/2.20 =
\qty{1.8}{\celsius}$ también.
% ledger: ar6ch7:tcr, ar6ch7:energy-budget
% values: tcr_energy 1.840 (anthropogenic.py numbers)

\section{Detección y atribución}\label{sec:b2:anthropogenic-climate-change:attribution}

\begin{definition}[Detección; atribución; huella climática]\label{def:b2:anthropogenic-climate-change:da}
La \emph{detección}\index{deteccion (del cambio climatico)@detección (del cambio climático)} de un cambio climático consiste en mostrar que el
clima ha cambiado en un sentido estadístico definido, más allá de lo que podría producir su variabilidad interna,
sin indicar la causa. La \emph{atribución}\index{atribucion@atribución} consiste en evaluar las
contribuciones de varias causas a un cambio detectado, con una probabilidad asignada. Una \emph{huella
climática}\index{huella climatica@huella climática} es la distribución en el espacio y el tiempo de la respuesta a un
factor, tal como la simulan los modelos climáticos, por la que se reconoce ese factor en las observaciones.
\end{definition}

Factores distintos dejan huellas distintas. Más gas de efecto invernadero calienta la troposfera y enfría
la estratosfera, porque una manta infrarroja más gruesa hace que las capas superiores radien más hacia el espacio;
un Sol más brillante calentaría ambas. Las observaciones muestran la primera distribución: la troposfera se ha calentado
desde al menos la década de 1950, y es \emph{prácticamente seguro} que la estratosfera se ha enfriado. Es
\emph{muy probable} que los gases de efecto invernadero humanos hayan sido el principal factor del calentamiento troposférico
desde 1979, y \emph{extremadamente probable} que el agotamiento de la capa de ozono haya sido el principal factor del
enfriamiento de la estratosfera inferior de 1979 a mediados de la década de 1990.
% ledger: ar6ch3:upper-air

\begin{method}[El método de la huella óptima]\label{met:b2:anthropogenic-climate-change:fingerprinting}
\begin{enumerate}
  \item Simúlese por separado la respuesta a cada factor: las huellas $x_1, \dots, x_k$, como
        vectores de valores en el espacio y el tiempo.
  \item Escríbanse las observaciones como $y = \sum_i \beta_i x_i + \varepsilon$, con $\varepsilon$ la
        variabilidad interna, cuya estadística se toma de largas simulaciones sin forzamiento.
  \item Estímense los factores de escala $\beta_i$ por mínimos cuadrados, ponderados por la inversa de la
        variabilidad para que las componentes más ruidosas cuenten menos, con sus intervalos de confianza.
  \item Un $\beta_i$ significativamente mayor que cero significa que la respuesta al factor $i$ está
        detectada; un intervalo que contiene 1 y excluye 0 significa que la respuesta simulada es
        coherente con las observaciones, y el cambio puede atribuirse en parte a ese factor.
\end{enumerate}
\end{method}
% ledger: ar6ch3:fingerprint

Aplicado a la temperatura de la superficie, el método da la \omterm{def:b2:anthropogenic-climate-change:da}{atribución} recordada al principio: de los
\qty{1.06}{\celsius} observados entre 1850--1900 y 2010--2019, la influencia humana causó, \emph{probablemente},
de 0.8 a \qty{1.3}{\celsius}, con una mejor estimación de \qty{1.07}{\celsius}; los factores naturales y la
variabilidad interna aportan como mucho unas pocas décimas de grado en uno u otro sentido
(\cref{exo:b2:anthropogenic-climate-change:9,exo:b2:anthropogenic-climate-change:10}).
% ledger: ipcc:attribution, ipcc:observed-2010-2019

\section{Escenarios y emisiones acumuladas}\label{sec:b2:anthropogenic-climate-change:scenarios}

\begin{definition}[Trayectoria Socioeconómica Compartida]\label{def:b2:anthropogenic-climate-change:ssp}
Una \emph{Trayectoria Socioeconómica Compartida}\index{Trayectoria Socioeconomica Compartida@Trayectoria Socioeconómica Compartida} (SSP, por sus siglas en inglés) es uno de cinco
relatos sobre cómo podrían evolucionar la población, las economías y la tecnología a lo largo del siglo. Los
escenarios de emisiones de la evaluación se denominan SSP$x$-$y$: $x$ designa la trayectoria, e $y$ el
forzamiento radiativo aproximado, en \unit{W.m^{-2}}, que el escenario alcanza en 2100.
\end{definition}
% ledger: ar6spm:ssp

Los cinco escenarios principales van de SSP1-1.9 a SSP5-8.5. La evaluación de la ciencia física los toma como
dados: ni pondera su probabilidad ni evalúa los relatos. Para cada uno, los modelos climáticos
dan un calentamiento proyectado para 2081--2100 respecto de 1850--1900, desde \qty{1.4}{\celsius} (muy probablemente de 1.0
a 1.8) con SSP1-1.9 hasta \qty{4.4}{\celsius} (de 3.3 a 5.7) con SSP5-8.5. Las proyecciones comparten una
ley sencilla.
% ledger: ar6spm:ssp, ipcc:proj-temp

\begin{definition}[Respuesta climática transitoria a las emisiones acumuladas de CO$_2$]\label{def:b2:anthropogenic-climate-change:tcre}
La \emph{respuesta climática transitoria a las emisiones acumuladas de \ce{CO2}}\index{respuesta climatica transitoria a las emisiones acumuladas de CO2@respuesta climática transitoria a las emisiones acumuladas de CO$_2$}
(TCRE, por sus siglas en inglés) es el calentamiento global de la superficie por unidad de dióxido
de carbono emitido en total desde el periodo preindustrial, en \unit{\celsius} por
\qty{1000}{Gt} de \ce{CO2}.
\end{definition}

\begin{proposition}[El calentamiento sigue a las emisiones acumuladas]\label{prop:b2:anthropogenic-climate-change:tcre}
Con buena aproximación, el calentamiento causado por el dióxido de carbono es proporcional al total emitido,
sea cual sea la trayectoria de las emisiones: $\Delta T \approx \text{TCRE} \times E$, con $E$ las emisiones
acumuladas.
\end{proposition}
\admitted
% ledger: ar6spm:tcre, ar6ch5:tcre-linear

La ley no es evidente. El forzamiento solo crece como el logaritmo de la concentración, así que cada
tonelada añadida calienta menos; pero el océano absorbe calor y carbono con una eficacia cada vez menor a medida que
se acumulan las emisiones, de modo que más parte de cada tonelada se queda en el aire y su calentamiento se retrasa menos. Los
dos efectos casi se compensan, y todos los estudios coinciden en que ese es el origen de la ley. La relación es
casi lineal (\emph{confianza alta}), y cada \qty{1000}{Gt} de dióxido de carbono causa, \emph{probablemente},
de 0.27 a \qty{0.63}{\celsius} de calentamiento, con una mejor estimación de \qty{0.45}{\celsius}, es decir, \qty{1.65}{\celsius}
por \qty{1000}{PgC} de carbono (\cref{exo:b2:anthropogenic-climate-change:4}). De 1850 a 2019 se emitieron
$\num{2390} \pm \qty{240}{Gt}$ de dióxido de carbono.
% ledger: ar6ch5:tcre-linear, ar6spm:tcre, ar6spm:budget

\begin{definition}[Presupuesto de carbono restante]\label{def:b2:anthropogenic-climate-change:budget}
El \emph{presupuesto de carbono restante}\index{presupuesto de carbono restante} para un límite de temperatura es el
dióxido de carbono acumulado que aún puede emitirse, a partir de una fecha dada, sin que el calentamiento global supere
ese límite con una probabilidad dada, teniendo en cuenta el calentamiento causado por los demás factores
humanos.
\end{definition}

\begin{method}[El presupuesto de carbono restante a partir de la TCRE]\label{met:b2:anthropogenic-climate-change:budget}
\begin{enumerate}
  \item Réstese del límite el calentamiento humano ya alcanzado, \qty{1.07}{\celsius} en
        2010--2019.
  \item Réstese el calentamiento que añadirán los demás factores hasta que las emisiones de dióxido de carbono
        alcancen el cero neto.
  \item Divídase por la TCRE: este es el presupuesto desde la mitad de 2010--2019.
  \item Réstense las emisiones desde entonces, y un margen para las retroalimentaciones lentas del sistema Tierra que
        los modelos omiten (el deshielo del permafrost es la mayor).
\end{enumerate}
\end{method}

La evaluación combina así cinco componentes: la TCRE; el calentamiento histórico; un
calentamiento adicional por los demás factores de unos 0.1 a \qty{0.2}{\celsius}; el calentamiento que puede
continuar después de que cesen las emisiones, evaluado en cero de media (\emph{probablemente} dentro de $\pm
\qty{0.19}{\celsius}$); y las retroalimentaciones del sistema Tierra, unas \qty{26}{Gt} de dióxido de carbono por
grado. Para \qty{1.5}{\celsius}, \qty{0.43}{\celsius} por encima de 2010--2019, el presupuesto restante desde el
comienzo de 2020 es de \qty{500}{Gt} de dióxido de carbono con una probabilidad del \qty{50}{\percent}, y de
\qty{1350}{Gt} para \qty{2}{\celsius}. Las decisiones sobre los demás factores desplazan estos valores en
\qty{220}{Gt} o más; el problema de fin de semana reconstruye el primero.
% ledger: ar6ch5:budget-parts, ar6spm:budget

\begin{dated}{2025}
% ledger: data:co2-mlo, data:gistemp (values read from the datasets as fetched 2026-10-22; offset computed)
La media anual del dióxido de carbono en Mauna Loa alcanzó \qty{427.35}{ppm} en 2025, un forzamiento de
$3.93\log_2(427.35/278) = \qty{2.44}{W.m^{-2}}$ según la ley logarítmica. En el registro de temperatura de la NASA,
2024 es el año más cálido, \qty{1.50}{\celsius} por encima de la media de 1880--1900, los primeros
años de ese registro; 2025 se situó en \qty{1.40}{\celsius}.
\end{dated}

\begin{omfigure}
\centering
\begin{tikzpicture}
\begin{axis}[omrecord, width=0.86\linewidth, height=5.8cm, xmin=0, xmax=4000, ymin=0, ymax=2.6,
  xlabel={emisiones acumuladas de \ce{CO2} desde 1850 (\unit{Gt})},
  ylabel={calentamiento desde 1850--1900 (\unit{\celsius})}, xtick={0,500,...,4000},
  /pgf/number format/1000 sep={\,}, xtick pos=bottom, ytick pos=left, clip=false,
  tick label style={font=\scriptsize}, label style={font=\small}]
% computed in figdata/bachelor-2/anthropogenic.py part tcre; ledger: ar6spm:tcre, ar6spm:budget
\addplot[name path=hi, draw=none, forget plot] table[x=E_GtCO2, y=dT_hi] {figdata/out/bachelor-2-anthropogenic-tcre.dat};
\addplot[name path=lo, draw=none, forget plot] table[x=E_GtCO2, y=dT_lo] {figdata/out/bachelor-2-anthropogenic-tcre.dat};
\addplot[omThm!18, forget plot] fill between[of=hi and lo];
\addplot[omThm, line width=1.3pt] table[x=E_GtCO2, y=dT_mid] {figdata/out/bachelor-2-anthropogenic-tcre.dat};
\draw[densely dashed, black!60] (axis cs:0,1.5) -- (axis cs:4000,1.5);
\draw[densely dashed, black!60] (axis cs:0,2.0) -- (axis cs:4000,2.0);
\node[anchor=north west, font=\scriptsize, text=black!70] at (axis cs:0,1.5) {\qty{1.5}{\celsius}};
\node[anchor=north west, font=\scriptsize, text=black!70] at (axis cs:0,2.0) {\qty{2}{\celsius}};
\fill[black] (axis cs:2390,1.07) circle (2.2pt);
\draw[black, line width=0.8pt] (axis cs:2150,1.07) -- (axis cs:2630,1.07);
\draw[black, line width=0.8pt] (axis cs:2390,0.8) -- (axis cs:2390,1.3);
\node[anchor=north west, font=\scriptsize, align=left] at (axis cs:2440,1.02) {1850--2019:\\\qty{2390}{Gt}, \qty{1.07}{\celsius}};
\draw[<->, omDef, thick] (axis cs:2390,1.62) -- (axis cs:2890,1.62);
\node[anchor=south, font=\scriptsize, text=omDef, align=center] at (axis cs:2640,1.64) {AR6: \qty{500}{Gt}};
\node[anchor=north west, font=\scriptsize, text=omThm, align=left] at (axis cs:150,2.5) {intervalo probable de la TCRE,\\de 0.27 a \qty{0.63}{\celsius} por \qty{1000}{Gt}};
\end{axis}
\end{tikzpicture}
\par\smallskip{\small Calentamiento frente a las emisiones acumuladas de dióxido de carbono para la TCRE evaluada,
\qty{0.45}{\celsius} por \qty{1000}{Gt} (línea), y su intervalo probable (sombreado), con el calentamiento humano
y las emisiones hasta la fecha (punto, con los intervalos probables como barras). La flecha marca el presupuesto restante para
\qty{1.5}{\celsius} con un \qty{50}{\percent} desde 2020: menor que una lectura de la línea sola,
porque los demás factores también calientan.}
\end{omfigure}


\section{El nivel del mar}\label{sec:b2:anthropogenic-climate-change:sea-level}

\begin{definition}[Balance del nivel del mar; subida termostérica del nivel del mar]\label{def:b2:anthropogenic-climate-change:sea-level}
El \emph{balance del nivel del mar}\index{balance del nivel del mar} es la comparación entre el cambio observado del
nivel medio global del mar y la suma de sus contribuciones estimadas: la dilatación del océano que se
calienta, la fusión de los glaciares y de los casquetes de hielo de Groenlandia y la Antártida, y el cambio del
agua almacenada en tierra. La \emph{subida termostérica del nivel del mar}\index{subida termosterica del nivel del mar@subida termostérica del nivel del mar} es la
parte debida a la dilatación térmica del agua de mar, a masa constante.
\end{definition}

\begin{proposition}[Dilatación térmica de una capa]\label{prop:b2:anthropogenic-climate-change:thermosteric}
Una capa de océano de espesor $H$, con coeficiente de dilatación térmica $\alpha$, que se calienta uniformemente
en $\Delta T$ a masa y superficie constantes, sube $\Delta h = \alpha\,\Delta T\,H$.
\end{proposition}
\begin{proof}
La masa $\rho H A$ de una columna de superficie $A$ es constante, así que $\Delta H/H = -\Delta\rho/\rho =
\alpha\,\Delta T$ a primer orden, con $\alpha = -(1/\rho)\,\partial\rho/\partial T$.
\end{proof}
% ledger: teos10:alpha-beta

En el océano real, $\alpha$ crece con la temperatura y la presión, y el calentamiento se concentra cerca de
la superficie. La evaluación resume el resultado como una eficiencia de dilatación de \qty{0.113}{m} por
\qty{e24}{J} de calor añadido al océano, o \qty{0.617}{m} por grado de calentamiento de todo el océano
(\emph{confianza media}). Sobre su profundidad media de \qty{3682}{m}, eso supone un $\alpha$ efectivo de
$0.617/3682 = \qty{1.68e-4}{\celsius^{-1}}$ (\cref{exo:b2:anthropogenic-climate-change:8}).
% ledger: ar6ch9:expansion, charette:ocean-volume

El balance, representado en la figura siguiente, se cierra dentro de sus incertidumbres. En 1971--2018, las contribuciones suman
\qty{2.00}{mm.yr^{-1}}, frente a una subida observada de \qty{2.33}{mm.yr^{-1}}; en 2006--2018
suman 3.61 frente a 3.69. La dilatación térmica fue el mayor término individual, la mitad de la suma en 1971--2018, pero el hielo
la está alcanzando: en 2006--2018, glaciares y casquetes de hielo aportaron en conjunto \qty{1.62}{mm.yr^{-1}}, más que
la dilatación. El océano sigue además calentándose en profundidad mucho después de que la superficie se estabilice, y los casquetes
de hielo responden a lo largo de siglos; el nivel del mar tiene comprometida una subida durante siglos o milenios y
seguirá alto durante miles de años (\emph{confianza alta}). En los próximos dos mil años, se
proyecta que suba de unos 2 a \qty{3}{m} si el calentamiento se mantiene en \qty{1.5}{\celsius}, de 2 a
\qty{6}{m} con \qty{2}{\celsius} y de 19 a \qty{22}{m} con \qty{5}{\celsius} (\emph{confianza baja}).
% ledger: ar6ch9:table95, ar6spm:sea-commit

\begin{omfigure}
\centering
\begin{tikzpicture}
\begin{axis}[ybar stacked, width=0.72\linewidth, height=6.2cm, ymin=0, ymax=4.3, bar width=26pt,
  xtick=data, symbolic x coords={1971--2018, 1993--2018, 2006--2018}, enlarge x limits=0.25,
  ylabel={ritmo de subida (\unit{mm.yr^{-1}})}, tick label style={font=\scriptsize},
  label style={font=\small}, ytick pos=left, xtick pos=bottom,
  legend style={at={(1.03,0.5)}, anchor=west, font=\scriptsize, draw=black!30},
  legend cell align=left, reverse legend]
% ledger: ar6ch9:table95
\addplot[fill=omDef!70, draw=omDef!90!black] coordinates {(1971--2018,1.01) (1993--2018,1.31) (2006--2018,1.39)};
\addlegendentry{dilatación térmica}
\addplot[fill=omDef!30, draw=omDef!70!black] coordinates {(1971--2018,0.44) (1993--2018,0.55) (2006--2018,0.62)};
\addlegendentry{glaciares}
\addplot[fill=omThm!55, draw=omThm!80!black] coordinates {(1971--2018,0.25) (1993--2018,0.43) (2006--2018,0.63)};
\addlegendentry{Groenlandia}
\addplot[fill=omThm!25, draw=omThm!70!black] coordinates {(1971--2018,0.14) (1993--2018,0.25) (2006--2018,0.37)};
\addlegendentry{Antártida}
\addplot[fill=black!20, draw=black!50] coordinates {(1971--2018,0.15) (1993--2018,0.31) (2006--2018,0.60)};
\addlegendentry{agua continental}
\node[fill=black, rotate=45, minimum size=5pt, inner sep=0pt] at (axis cs:1971--2018,2.33) {};
\node[fill=black, rotate=45, minimum size=5pt, inner sep=0pt] at (axis cs:1993--2018,3.25) {};
\node[fill=black, rotate=45, minimum size=5pt, inner sep=0pt] at (axis cs:2006--2018,3.69) {};
\end{axis}
\end{tikzpicture}
\par\smallskip{\small El \omterm{def:b2:anthropogenic-climate-change:sea-level}{balance del nivel del mar}: contribuciones observadas a la subida del nivel medio global
del mar, como ritmos en tres periodos que terminan en 2018, apiladas, y el ritmo observado (rombos). La
diferencia entre la pila y el rombo es la parte que el balance no cierra. Valores: IPCC, Sexto Informe de
Evaluación, Grupo de Trabajo~I, tabla 9.5.}
\end{omfigure}

\section*{Ejercicios}\label{sec:b2:anthropogenic-climate-change:exercises}

Datos de los ejercicios, salvo aviso contrario: forzamiento de una duplicación del dióxido de carbono $F_{2\times} =
\qty{3.93}{W.m^{-2}}$ desde $C_0 = \qty{278}{ppm}$; constante de Stefan--Boltzmann $\sigma =
\qty{5.670e-8}{W.m^{-2}.K^{-4}}$; radio de la Tierra \qty{6371}{km}; un año,
\qty{3.156e7}{s}; masas molares del carbono y del dióxido de carbono, 12.011 y \qty{44.009}{g.mol^{-1}};
las retroalimentaciones y el \omterm{def:b2:anthropogenic-climate-change:sea-level}{balance del nivel del mar} del capítulo.
% ledger: ipcc:erf-2xco2, const:sigma, earth:Rmean, ciaaw:abridged-2024, ar6ch7:table710, ar6ch9:table95

\begin{exercise}[$\star$]\label{exo:b2:anthropogenic-climate-change:1}
El forzamiento de una duplicación del dióxido de carbono es de \qty{3.75}{W.m^{-2}} una vez ajustada la estratosfera;
la troposfera se ajusta después en $-0.60$ (temperatura), +0.22 (vapor de agua), +0.45 (nubes) y
+0.11 (superficie terrestre) \unit{W.m^{-2}}. (a) Hállese el ERF. (b) ¿En qué porcentaje supera a los
\qty{3.71}{W.m^{-2}} de la evaluación anterior? (c) ¿Qué ajuste es negativo, y por qué una troposfera más caliente
reduce el forzamiento?
\end{exercise}
% ledger: ar6ch7:2xco2

\begin{exercise}[$\star$]\label{exo:b2:anthropogenic-climate-change:2}
El sistema climático ganó \qty{434.9}{ZJ} de 1971 a 2018 y \qty{152.4}{ZJ} de 2006 a 2018,
contados entre las mitades del primer y del último año. (a) Hállese la superficie de la Tierra.
(b) Conviértase cada ganancia en un desequilibrio medio en \unit{W.m^{-2}}. (c) El océano se llevó el
\qty{91}{\percent}: ¿cuántos zettajulios en 1971--2018?
\end{exercise}
% ledger: ar6ch7:inventory, earth:Rmean

\begin{exercise}[$\star$]\label{exo:b2:anthropogenic-climate-change:3}
(a) Súmense los valores centrales de las retroalimentaciones de la figura, y compárese con la neta. (b) Hállense
la ECS con el valor neto y el calentamiento con la \omterm{def:b2:anthropogenic-climate-change:feedbacks}{respuesta de Planck} sola. (c) ¿En qué factor amplifican
las demás retroalimentaciones el calentamiento de Planck?
\end{exercise}

\begin{exercise}[$\star$]\label{exo:b2:anthropogenic-climate-change:4}
(a) Conviértase una TCRE de \qty{0.45}{\celsius} por \qty{1000}{Gt} de dióxido de carbono en grados por
\qty{1000}{PgC} de carbono, y hágase lo mismo con su intervalo probable, de 0.27 a 0.63. (b) Conviértanse las
\qty{2390}{Gt} de dióxido de carbono emitidas de 1850 a 2019 en petagramos de carbono. (c) ¿Qué
calentamiento da la TCRE para estas emisiones?
\end{exercise}
% ledger: ar6spm:tcre, ar6spm:budget

\begin{exercise}[$\star\star$]\label{exo:b2:anthropogenic-climate-change:5}
(a) La superficie emite \qty{398}{W.m^{-2}} y el efecto invernadero retiene \qty{159}{W.m^{-2}} de
ellos. Hállense $g$ y la temperatura de la superficie de un cuerpo negro que emite \qty{398}{W.m^{-2}}. (b) Hállese la
\omterm{def:b2:anthropogenic-climate-change:feedbacks}{respuesta de Planck} a \qty{288}{K}, y a \qty{298}{K} con el mismo $g$. (c) ¿Haría la \omterm{def:b2:anthropogenic-climate-change:feedbacks}{respuesta de Planck}
por sí sola más o menos sensible a una Tierra más caliente?
\end{exercise}
% ledger: ipcc:budget-surface, ar6ch7:planck

\begin{exercise}[$\star\star$]\label{exo:b2:anthropogenic-climate-change:6}
Los satélites dan una \omterm{def:b2:anthropogenic-climate-change:feedbacks}{retroalimentación del vapor de agua} de $1.85 \pm 0.32$ y los estudios una retroalimentación del gradiente vertical de
$-0.50 \pm 0.20$ \unit{W.m^{-2}.\celsius^{-1}} (una desviación típica). (a) Hállese su suma.
(b) Si fueran independientes, ¿cuál sería la desviación típica de la suma? (c) Los modelos
dan a la suma una desviación típica de solo \num{0.10}. ¿Qué dice eso de los dos errores?
\end{exercise}
% ledger: ar6ch7:wv-lr

\begin{exercise}[$\star\star$]\label{exo:b2:anthropogenic-climate-change:7}
De 1850--1900 a 2006--2019, $\Delta N = \qty{0.59}{W.m^{-2}}$, $\Delta T = \qty{1.03}{\celsius}$
y $\Delta F = \qty{2.20}{W.m^{-2}}$. (a) Hállense $\lambda$ y la sensibilidad efectiva. (b) Se espera que el
parámetro de retroalimentación se vuelva menos negativo en $\lambda' = \qty{0.5}{W.m^{-2}.\celsius^{-1}}$
a medida que evolucione la distribución del calentamiento. Hállese la sensibilidad $-F_{2\times}/(\lambda + \lambda')$.
(c) Rehágase (b) con $\lambda' = 0.1$.
\end{exercise}
% ledger: ar6ch7:energy-budget, ar6ch7:pattern, ar6ch7:pattern-ecs

\begin{exercise}[$\star\star$]\label{exo:b2:anthropogenic-climate-change:8}
(a) A partir de los \qty{0.617}{m} de subida por grado de calentamiento de todo el océano, de profundidad media
\qty{3682}{m}, hállese el coeficiente de dilatación efectivo. (b) El océano ganó \qty{396}{ZJ} en
1971--2018. Con una eficiencia de dilatación de \qty{0.113}{m} por \qty{e24}{J}, hállese la subida termostérica
y compárese con los \qty{47.5}{mm} del balance. (c) Una capa de \qty{700}{m} se calienta
\qty{0.2}{\celsius} con $\alpha = \qty{2e-4}{\celsius^{-1}}$: ¿cuánto hace subir el nivel del mar?
\end{exercise}
% ledger: ar6ch9:expansion, charette:ocean-volume, ar6ch7:inventory, ar6ch9:table95

\begin{exercise}[$\star\star$]\label{exo:b2:anthropogenic-climate-change:9}
Una huella de un modelo, en cuatro décadas, es $x = (0.10, 0.20, 0.30, 0.40)$ \unit{\celsius}; las
observaciones son $y = (0.12, 0.18, 0.33, 0.41)$ \unit{\celsius}. (a) Hállese el factor de escala de mínimos
cuadrados $\beta$ de $y = \beta x + \varepsilon$. (b) Su error típico, estimado a partir de simulaciones sin
forzamiento, es \num{0.08}. ¿Está detectada la respuesta? ¿Es coherente con el modelo? (c) ¿Qué
significaría un $\beta$ de 0.5 con el mismo error?
\end{exercise}

\begin{exercise}[$\star\star\star$]\label{exo:b2:anthropogenic-climate-change:10}
Dos huellas, la de los gases de efecto invernadero $g = (0.2, 0.4, 0.6, 0.8)$ y la de los aerosoles $a = (-0.1, -0.2,
-0.2, -0.3)$, ambas en \unit{\celsius} a lo largo de cuatro décadas, se ajustan a las observaciones $y = (0.1,
0.2, 0.4, 0.5)$. (a) Escríbanse las ecuaciones normales del ajuste por mínimos cuadrados $y = \beta_g g + \beta_a a
+ \varepsilon$ y resuélvanse. (b) Hállese el calentamiento atribuido a cada factor en la última década.
(c) ¿Por qué crecen las incertidumbres de $\beta_g$ y $\beta_a$ cuando las dos huellas se parecen?
\end{exercise}

\begin{exercise}[$\star\star\star$]\label{exo:b2:anthropogenic-climate-change:11}
Úsese el método del \omterm{def:b2:anthropogenic-climate-change:budget}{presupuesto de carbono restante} con la TCRE de \qty{0.45}{\celsius} por
\qty{1000}{Gt}, las \qty{210}{Gt} emitidas de 2015 a 2019 y \qty{26}{Gt} por grado para las
retroalimentaciones del sistema Tierra. (a) Hállese el presupuesto para \qty{2}{\celsius}, \qty{0.93}{\celsius} por encima de
2010--2019, si los demás factores añaden \qty{0.1}{\celsius}. (b) ¿Qué calentamiento adicional de los demás
factores lo iguala a las \qty{1350}{Gt} evaluadas? ¿Está dentro del intervalo evaluado? (c) Demuéstrese
que cada \qty{0.1}{\celsius} de calentamiento no debido al CO$_2$ cambia cualquier presupuesto en la misma cantidad.
\end{exercise}

\begin{exercise}[$\star\star\star$]\label{exo:b2:anthropogenic-climate-change:12}
La retroalimentación neta tiene un intervalo muy probable de $-1.81$ a \qty{-0.51}{W.m^{-2}.\celsius^{-1}}, en torno a
$-1.16$. (a) Hállese la ECS en los dos extremos y en el centro. (b) Demuéstrese que un intervalo de $\lambda$
simétrico respecto de su centro da un intervalo de ECS que no lo es, e indíquese qué lado es más largo. (c) ¿Por qué
les cuesta más a las evaluaciones descartar una ECS alta que una baja?
\end{exercise}
% ledger: ar6ch7:table710

\begin{problem}[{Problema de fin de semana --- La lectura del balance}]\label{pb:b2:anthropogenic-climate-change:1}
Este problema sigue la energía y el carbono de la evaluación desde el forzamiento hasta el presupuesto
restante para \qty{1.5}{\celsius}. Úsense los datos de los ejercicios, y: ERF 1750--2019 (\unit{W.m^{-2}}):
dióxido de carbono 2.16, metano 0.54, óxido nitroso 0.21, gases halogenados 0.41, ozono 0.47,
vapor de agua estratosférico 0.05, aerosol--radiación $-0.22$, aerosol--nube $-0.84$, uso del suelo
$-0.20$, carbono negro sobre la nieve 0.08, estelas de condensación 0.06; dióxido de carbono \qty{410}{ppm} en 2019 y
\qty{427.35}{ppm} en 2025; calentamiento del sistema \qty{0.79}{W.m^{-2}} en 2006--2018; los cambios del balance de energía
$\Delta N = 0.59$, $\Delta F = \qty{2.20}{W.m^{-2}}$, $\Delta T = \qty{1.03}{\celsius}$; TCRE
\qty{0.45}{\celsius} por \qty{1000}{Gt} de dióxido de carbono (probable de 0.27 a 0.63); \qty{2390}{Gt}
emitidas en 1850--2019; calentamiento humano \qty{1.07}{\celsius} en 2010--2019; \qty{210}{Gt} emitidas de
2015 a 2019; calentamiento adicional no debido al CO$_2$ \qty{0.1}{\celsius}, el extremo inferior de los 0.1
a \qty{0.2}{\celsius} evaluados; retroalimentaciones del sistema Tierra \qty{26}{Gt} por grado. Redondéese a tres cifras
significativas.
% ledger: ipcc:erf-agents, ipcc:erf-2xco2, data:co2-mlo, ar6ch7:inventory, ar6ch7:energy-budget, ar6spm:tcre, ar6spm:budget, ar6ch5:budget-parts

\noindent\textbf{Parte I --- El forzamiento.}
\begin{enumerate}
  \item Súmese el forzamiento de los cuatro gases de efecto invernadero bien mezclados.
  \item Súmese el forzamiento de los aerosoles.
  \item Súmense todos los términos humanos. Compárese con los \qty{2.72}{W.m^{-2}} evaluados.
  \item Compruébese el forzamiento del dióxido de carbono en 2019 con la ley logarítmica.
  \item Hállese el forzamiento del dióxido de carbono en 2025 con la misma ley.
  \item ¿Qué término tiene la incertidumbre más amplia, y por qué domina la del total?
\end{enumerate}

\noindent\textbf{Parte II --- El balance de energía.}
\begin{enumerate}[resume]
  \item ¿Cuánta energía ganó el sistema climático en 2006--2018, en zettajulios?
  \item ¿Cuánta de ella fue al océano?
  \item Hállese el parámetro de retroalimentación a partir de los cambios del balance de energía.
  \item Hállese la \omterm{def:b2:anthropogenic-climate-change:ecs}{sensibilidad climática de equilibrio} efectiva.
  \item Hállese la \omterm{def:b2:anthropogenic-climate-change:tcr}{respuesta climática transitoria} a partir de los mismos datos.
  \item Hállese la ECS a partir de la retroalimentación neta de la figura, y compárense los dos valores.
  \item Dese la razón física por la que la estimación histórica es la más baja.
\end{enumerate}

\noindent\textbf{Parte III --- El calentamiento por tonelada.}
\begin{enumerate}[resume]
  \item Conviértase la TCRE a grados por \qty{1000}{PgC} de carbono.
  \item Conviértase su intervalo probable.
  \item Conviértanse las emisiones de 1850--2019 en petagramos de carbono.
  \item ¿Qué calentamiento da la TCRE para las emisiones de 1850--2019?
  \item Compárese con el calentamiento humano de \qty{1.07}{\celsius}. ¿Qué dice la comparación del
        efecto neto de los demás factores humanos?
\end{enumerate}

\noindent\textbf{Parte IV --- El presupuesto restante para \qty{1.5}{\celsius}.}
\begin{enumerate}[resume]
  \item ¿Qué calentamiento queda por encima de 2010--2019?
  \item Quítese el calentamiento no debido al CO$_2$ y divídase por la TCRE: ¿qué presupuesto queda desde la mitad
        de 2010--2019?
  \item Quítense las emisiones de 2015--2019.
  \item Quítense las retroalimentaciones del sistema Tierra sobre el calentamiento restante.
  \item ¿Cuánto cambia el presupuesto si los demás factores añaden \qty{0.2}{\celsius} en su lugar?
  \item En resumen: ¿cuál es el \omterm{def:b2:anthropogenic-climate-change:budget}{presupuesto de carbono restante} para \qty{1.5}{\celsius}, desde el comienzo de
        2020, calculado a partir de la TCRE? Compárese con las \qty{500}{Gt} evaluadas.
\end{enumerate}
\end{problem}
